% 240 words as plain text (the formula as two words, each number with its percent sign as one);
% 249 if every percent sign counts as a word. The journal's limit is 250.
Soft sensors are recalibrated against laboratory or analyzer measurements that arrive late.
Updated with such labels, the single-state recursion of an adaptive conformal interval becomes an
integral loop with dead time: for a delay of $\theta$ updates its linearisation is stable only for
step sizes below $2\sin(\pi/(4\theta+2))$, and on serially correlated residuals it can issue
uninformative intervals while its marginal coverage stays nominal. We test two remedies in a
pre-registered study on 15 Tennessee Eastman runs at three operating points and the debutanizer
data. With the Tennessee Eastman analyzer four updates late, steps tuned without the delay left
6.9\,\% of intervals uninformative, 6.0 percentage points more than without delay, at coverage
0.897. Replaying historical residuals through the delayed loop to tune the step held a 2\,\% budget of uninformative intervals (0.30\,\%). Predicting the persistent
part of the residual and calibrating only the prediction error narrowed the intervals by
18\,\% (95\,\% CI 8 to 28\,\%) at the same budget (0.18\,\%), made misses less persistent and
allowed larger steps; its width was within 1.5\,\% of the best of four conformal PID control
and SPCI configurations. Predicted from the calibration residuals' autocorrelation, the
width ratio erred by $-4.9$\,\% on average at four updates and was conservative at every delay
checked. The delay did not raise the uninformative share on the debutanizer at its primary dead
time, so the diagnosis is restricted to the simulated plant, where compensated split conformal
matched the loop.
